% !TEX root = ../matan.tex

\section{Теорема Дини.}

\begin{Theorem}{Теорема Дини}{}
	Пусть $I = [a, b]$, функции $f_n \colon I \to \RR$ непрерывны, $f_n(x) \ge 0$ для всех $x \in I$, и при каждом $x \in I$ последовательность $\{f_n(x)\}$ монотонно убывает к нулю:
	\[
	f_n(x) \ge f_{n+1}(x), \qquad f_n(x) \to 0 \quad (n \to \infty).
	\]
	Тогда $f_n \to 0$ равномерно на $I$.
\end{Theorem}

\begin{proof}
	Зафиксируем $\eps > 0$. Для каждой точки $x \in I$ из поточечной сходимости $f_n(x) \to 0$ найдётся номер $N(x)$ такой, что
	\[
	f_{N(x)}(x) < \eps.
	\]
	Функция $f_{N(x)}$ непрерывна в точке $x$, поэтому существует окрестность $U_x$ точки $x$, в которой
	\[
	f_{N(x)}(y) < 2\eps \qquad \forall y \in U_x \cap I.
	\]

	Семейство открытых множеств $\{U_x\}_{x \in I}$ покрывает компакт $I$, поэтому из него можно выделить конечное подпокрытие
	\[
	I \subset U_{x_1} \cup \cdots \cup U_{x_M}.
	\]
	Положим $\widetilde{N} = \max\{N(x_1), \ldots, N(x_M)\}$.

	Возьмём произвольные $y \in I$ и $n > \widetilde{N}$. Точка $y$ принадлежит некоторому $U_{x_j}$, поэтому $f_{N(x_j)}(y) < 2\eps$. По монотонности (последовательность $\{f_k(y)\}$ убывает, и $n > \widetilde{N} \ge N(x_j)$)
	\[
	0 \le f_n(y) \le f_{N(x_j)}(y) < 2\eps.
	\]
	Так как оценка верна для всех $y \in I$, имеем $\sup_{y \in I} f_n(y) \le 2\eps$ при всех $n > \widetilde{N}$. Следовательно, $f_n \to 0$ равномерно на $I$.
\end{proof}

\begin{Remark}{}{}
	Все три условия существенны: компактность $I$ (без неё конечного подпокрытия может не быть), непрерывность $f_n$ (и предельной функции), монотонность. Например, на некомпактном или при немонотонной сходимости заключение, вообще говоря, неверно.
\end{Remark}

\begin{Corollary}{Признак Дини для рядов}{}
	Пусть $I = [a, b]$, функции $f_k \colon I \to \RR$ непрерывны и $f_k(x) \ge 0$ для всех $x \in I$. Если ряд $\sum_{k=1}^{\infty} f_k(x)$ сходится поточечно на $I$ к \textbf{непрерывной} функции $f$, то он сходится равномерно на $I$.
\end{Corollary}

\begin{proof}
	Обозначим частичные суммы $S_n(x) = \sum_{k=1}^{n} f_k(x)$ и остатки $r_n(x) = f(x) - S_n(x) = \sum_{k=n+1}^{\infty} f_k(x)$. Так как $f_k \ge 0$, суммы $S_n$ не убывают по $n$, поэтому остатки $r_n$ не возрастают: $r_n(x) \ge r_{n+1}(x) \ge 0$. Каждая $r_n = f - S_n$ непрерывна (как разность непрерывных), и $r_n(x) \to 0$ поточечно. По теореме Дини $r_n \to 0$ равномерно, то есть
	\[
	\sup_{x \in I} |f(x) - S_n(x)| \to 0,
	\]
	что и означает равномерную сходимость ряда $\sum f_k$ к $f$.
\end{proof}
